\section{The semiregularity target}\label{sec:semireg}

Markman's construction at $n=2$ proceeds by deforming a sheaf across a family
and controlling the obstruction by a semiregularity argument in the sense of
Bloch and of Buchweitz and Flenner \cite{Blo72,BF03}. It is sometimes said
that this route is blocked for $n\ge4$ by a dimension count, the target of the
semiregularity map being too large. We compute the target and show that the
count says nothing of the kind. The section is included because the belief
that the semiregularity route is numerically blocked, combined with the belief
that the secant route is numerically blocked, leaves the impression that the
difficulty at $n\ge4$ is arithmetic. It is not: it is
\Cref{thm:divisor,thm:character}.

\begin{lemma}[Hodge numbers of an abelian variety]\label{lem:hodgenumbers}
For an abelian variety $A$ of dimension $g$ one has
$h^{p,q}(A)=\binom{g}{p}\binom{g}{q}$.
\end{lemma}

\begin{proof}
$H^{p,q}(A)=\bigwedge^{p}H^{1,0}\otimes\bigwedge^{q}H^{0,1}$ by
\eqref{eq:wedge} and the compatibility of the Hodge decomposition with the
exterior algebra structure, and $\dim H^{1,0}=\dim H^{0,1}=g$.
\end{proof}

\begin{proposition}[The semiregularity target]\label{prop:target}
Let $A$ be an abelian variety of dimension $g=2n$. The target of the
semiregularity map, namely $\bigoplus_{q}H^{q+2}(A,\Omega^{q}_{A})$, has
dimension
\[
  \sum_{q\ge0}h^{q,q+2}(A)=\binom{4n}{2n-2}.
\]
\end{proposition}

\begin{proof}
By \Cref{lem:hodgenumbers} with $g=2n$,
\[
  \sum_{q\ge0}h^{q,q+2}
  =\sum_{q\ge0}\binom{2n}{q}\binom{2n}{q+2}
  =\sum_{q\ge0}\binom{2n}{q}\binom{2n}{2n-q-2}
  =\binom{4n}{2n-2},
\]
the last step being the Vandermonde convolution.
\end{proof}

\begin{corollary}[The count does not bind]\label{cor:countnotbind}
The target dimension $\binom{4n}{2n-2}$ takes the values
$28,495,8008,125970,1961256,30421755$ for $n=2,\dots,7$, growing by a factor
exceeding ten at each step. Any source whose dimension grows polynomially in
$n$, for instance the quadratic quantity $\binom{2n}{2}=n(2n-1)$ which takes
the values $6,15,28,45,66,91$ over the same range, is eventually negligible
against it. In particular a comparison of the two dimensions cannot obstruct
the existence of a semiregular sheaf for large $n$; the inequality it produces
points the other way.
\end{corollary}

\begin{remark}[Consistency at $n=2$]\label{rem:markmanconsistent}
At $n=2$ the target is $28$ and the sheaf produced by Markman has
$\dim\Ext^{2}=12$, comfortably inside it, which is consistent with the
construction succeeding there. Nothing in the numbers changes character
between $n=2$ and $n=4$. Whatever prevents the construction from being carried
out for larger $n$, it is not this comparison.
\end{remark}

\begin{remark}[What semiregularity would still need]\label{rem:semiregneed}
We emphasise that \Cref{cor:countnotbind} is not an argument that the
semiregularity route works. It removes one purported obstruction and leaves
the real requirement in place, which is the construction of a sheaf on $A$
whose Chern character has a nonzero component along the Weil line. By
\Cref{thm:character} such a sheaf cannot be $K$-equivariant with the norm
character, so it cannot be produced by any functorial construction from
$K$-invariant data on $\Av$; Markman's sheaf at $n=2$ evades this because it
is not so produced. That is the sense in which the problem for $n\ge3$ is a
problem about breaking equivariance, and not about dimensions.
\end{remark}

\subsection{The annihilator of the Weil line, and objects of rank zero}
\label{ssec:rankzero}

\Cref{rem:semiregneed} says that the remaining requirement is a sheaf whose
Chern character has a component along the Weil line. The most economical way
to arrange that is to make the Chern character lie \emph{entirely} on the
Weil line: to take an object of rank zero, with every $\ch_{k}$ vanishing
except $\ch_{n}$, which is what the construction of \Cref{sec:construction}
does. We show here that this economy is fatal to semiregularity, for a
reason that has nothing to do with how the object is built and everything to
do with where its Chern character sits.%
\footnote{We recall what the map is and what it is for, since the rest of
this subsection turns on it. For a coherent sheaf $\cE$ on a smooth
projective variety $A$ the semiregularity map is defined through the Atiyah
class $\mathrm{at}(\cE)\in\Ext^{1}(\cE,\cE\otimes\Omega^{1}_{A})$: a class
$\xi\in\Ext^{2}(\cE,\cE)$ is sent to the trace of $\xi\cdot\exp(\mathrm{at}(\cE))$,
and the component of that trace in $H^{q}(A,\Omega^{q+2}_{A})=H^{q,q+2}(A)$ is the
$q$-th coordinate of $\sigma_{\cE}(\xi)$. Bloch introduced the map for local
complete intersections \cite{Blo72}, and Buchweitz and Flenner extended it to
arbitrary coherent sheaves and perfect complexes \cite{BF03}. Its use is this. If
$A$ moves in a family $\cA\to S$, the obstruction to extending $\cE$ to first
order along $v\in T_{s}S$ is a class $\mathrm{ob}_{v}(\cE)\in\Ext^{2}(\cE,\cE)$,
and the theorem of Buchweitz and Flenner identifies $\sigma_{\cE}(\mathrm{ob}_{v}(\cE))$
with the derivative of $\ch(\cE)$ in the direction $v$, projected off the
diagonal of the Hodge decomposition. So if $\ch(\cE)$ stays of Hodge type along
$S$ then $\sigma_{\cE}(\mathrm{ob}_{v}(\cE))=0$, and if $\sigma_{\cE}$ is also
injective then $\mathrm{ob}_{v}(\cE)=0$ and $\cE$ deforms. Injectivity is what
converts a statement about the Chern character into a statement about the
sheaf, and it is the one hypothesis in the argument that sees more than
$\ch(\cE)$. Two special cases are all that is needed below: for a line bundle
$L$, $\Ext^{2}(L,L)=H^{2}(\cO_{A})=H^{0,2}(A)$ and $\sigma_{L}$ is multiplication
by $e^{c_{1}(L)}$; and for a direct sum, $\Ext^{2}$ is the sum of the
$\Ext^{2}(\cE_{i},\cE_{j})$ and $\sigma$ acts on each diagonal block by
multiplication by the Chern character of that summand. Both follow at once from
the definition by the trace, and both are what items (XIII) and (XIV) of
\Cref{sec:verification} implement.}

\begin{lemma}[The annihilator of the Weil line in $H^{0,2}$]
\label{lem:annihilator}
Let $A$ be of $(K,-1,n)$-Weil type and let $\omega\in\HW(A)\otimes\CC$ be
any nonzero element of the Weil line. Then the multiplication map
\[
  H^{0,2}(A)\longrightarrow H^{n,n+2}(A),\qquad z\longmapsto z\cdot\omega,
\]
has kernel of dimension exactly $n^{2}$. Its kernel is spanned by the
products $u\wedge v$ with $u\in V_{+}^{0,1}$ and $v\in V_{-}^{0,1}$, and it
is the same subspace for every $\omega$ on the line.
\end{lemma}

\begin{proof}
Write $\omega=a\,\alpha_{+}+b\,\alpha_{-}$ with $\alpha_{\pm}$ generating
$\bigwedge^{2n}V_{\pm}$ and $(a,b)\neq(0,0)$. Since
$H^{0,1}=V_{+}^{0,1}\oplus V_{-}^{0,1}$ with each summand of dimension $n$,
the space $H^{0,2}=\bigwedge^{2}H^{0,1}$ decomposes as
\[
  H^{0,2}
  =\textstyle\bigwedge^{2}V_{+}^{0,1}
  \;\oplus\;\bigl(V_{+}^{0,1}\otimes V_{-}^{0,1}\bigr)
  \;\oplus\;\bigwedge^{2}V_{-}^{0,1},
\]
of dimensions $\binom{n}{2}$, $n^{2}$ and $\binom{n}{2}$.
A mixed product $u\wedge v$ with $u\in V_{+}^{0,1}$, $v\in V_{-}^{0,1}$
satisfies $u\wedge v\wedge\alpha_{+}=0$ because $u$ already occurs among
the $2n$ factors of $\alpha_{+}$, and $u\wedge v\wedge\alpha_{-}=0$ for the
same reason with $v$. So the middle summand lies in the kernel.

Conversely let $z=z_{++}+z_{+-}+z_{--}$ be in the kernel. The product
$z\cdot\omega$ has the component $a\,z_{--}\wedge\alpha_{+}$ in
$\bigwedge^{2}V_{-}^{0,1}\otimes\bigwedge^{2n}V_{+}$ and the component
$b\,z_{++}\wedge\alpha_{-}$ in
$\bigwedge^{2}V_{+}^{0,1}\otimes\bigwedge^{2n}V_{-}$, and these two summands
of $\bigwedge^{2n+2}V_{\CC}$ are independent of each other and of everything
else in $z\cdot\omega$. Each of the maps $z_{--}\mapsto z_{--}\wedge\alpha_{+}$
and $z_{++}\mapsto z_{++}\wedge\alpha_{-}$ is injective, since wedging a
nonzero element of $\bigwedge^{2}V_{-}^{0,1}$ with a basis of $V_{+}$
produces a nonzero element of $\bigwedge^{2n+2}(V_{+}\oplus V_{-})$. Hence
if $a\neq0$ then $z_{--}=0$, and if $b\neq0$ then $z_{++}=0$. When both
$a$ and $b$ are nonzero the kernel is exactly the middle summand. When one
of them vanishes, say $b=0$, the class $\omega=a\,\alpha_{+}$ is not real
and lies in the complexified Weil line; the kernel is then
$V_{+}^{0,1}\otimes V_{-}^{0,1}\oplus\bigwedge^{2}V_{+}^{0,1}$, of dimension
$n^{2}+\binom{n}{2}$. For the rational generators of \Cref{thm:explicitgens},
where $a$ and $b$ are both nonzero, the kernel has dimension exactly $n^{2}$.
\end{proof}

The dimension $n^{2}$ was also computed directly in the exterior algebra,
modulo a prime, for both rational generators of the Weil line at
$n=2,3,4$, and agrees (\Cref{sec:verification}, item (XIV)).

\begin{theorem}[Objects of rank zero are not semiregular]
\label{thm:rankzero}
Let $A$ be of $(K,-1,n)$-Weil type. Let $\cE_{1}$ and $\cE_{2}$ be coherent
sheaves on $A$ such that
\[
  \ch(\cE_{1})-\ch(\cE_{2})=c\,\omega,\qquad c\neq0,\quad \omega\in\HW(A),
\]
and such that $\Ext^{2}(\cE_{1},\cE_{2})=\Ext^{2}(\cE_{2},\cE_{1})=0$. Put
$\cE=\cE_{1}\oplus\cE_{2}$. Then the semiregularity map
\[
  \sigma_{\cE}\colon\ \Ext^{2}(\cE,\cE)\longrightarrow\bigoplus_{q\ge0}H^{q,q+2}(A)
\]
has kernel of dimension at least $n^{2}$. In particular $\cE$ is not
semiregular, for any $n\ge1$.
\end{theorem}

\begin{proof}
The vanishing of the off-diagonal $\Ext^{2}$ groups gives
$\Ext^{2}(\cE,\cE)=\Ext^{2}(\cE_{1},\cE_{1})\oplus\Ext^{2}(\cE_{2},\cE_{2})$,
and each summand contains $H^{2}(\cO_{A})\cdot\id_{\cE_{i}}$ through the
$\cO_{A}$-module structure, the map $z\mapsto z\cdot\id_{\cE_{i}}$ being
injective for $\cE_{i}\neq0$.%
\footnote{The map $H^{2}(\cO_{A})\to\Ext^{2}(\cE,\cE)$, $z\mapsto z\cdot\id_{\cE}$,
is the edge map of the local-to-global spectral sequence, and when $\rk\cE\ne0$ it
is split by the trace up to the factor $\rk\cE$. For a direct sum of line bundles
every summand has rank one, so the map is injective on each block with no
hypothesis, and the element $\xi_{z}$ of the proof, which is $z\cdot\id$ on the
first block and $-z\cdot\id$ on the second, is nonzero for every nonzero $z$.
This is the only place the direct-sum structure is used: it is what allows an
element of $\Ext^{2}$ whose image under $\sigma$ is
$z\cdot(\ch\cE_{1}-\ch\cE_{2})$ rather than $z\cdot(\ch\cE_{1}+\ch\cE_{2})$.
For a single indecomposable sheaf of rank zero the same conclusion holds by the
first sentence of the proof applied to $z\cdot\id_{\cE}$, provided that element
is nonzero, which is the case whenever the restriction of $H^{0,2}(A)$ to the
support of $\cE$ is injective; it fails for a skyscraper sheaf, whose support is
a point, and the theorem is stated for direct sums to keep that hypothesis in
view.} For $z\in H^{2}(\cO_{A})=H^{0,2}(A)$ consider
the element $\xi_{z}=z\cdot\id_{\cE_{1}}-z\cdot\id_{\cE_{2}}$. The
semiregularity map is compatible with direct sums and on the summand
$H^{0,2}\cdot\id_{\cE_{i}}$ it is multiplication by the Chern character of
$\cE_{i}$, so
\[
  \sigma_{\cE}(\xi_{z})=z\cdot\ch(\cE_{1})-z\cdot\ch(\cE_{2})
  =c\,z\cdot\omega .
\]
By \Cref{lem:annihilator} this vanishes for every $z$ in an $n^{2}$-dimensional
subspace of $H^{0,2}$, and the corresponding $\xi_{z}$ are nonzero and
linearly independent. Hence $\dim\ker\sigma_{\cE}\ge n^{2}$.
\end{proof}

\begin{corollary}[The construction of \Cref{sec:construction}]
\label{cor:pteobstructed}
Let $L_{1},\dots,L_{m}$ be line bundles on $A$ with
$\sum_{i}c_{1}(L_{i})=0$ such that the virtual class
$\sum_{i}\bigl([L_{i}]-[L_{i}^{-1}]\bigr)$ has Chern character on the Weil
line and every difference $c_{1}(L_{i})\pm c_{1}(L_{j})$, $i\neq j$, is
nondegenerate of index different from $2$. Then the sheaf
$\bigoplus_{i}(L_{i}\oplus L_{i}^{-1})$ has semiregularity kernel of
dimension at least $n^{2}$. At $n=3$ the explicit object of
\Cref{sec:construction} has kernel of dimension exactly
$15=n^{2}+n(n-1)$, of which the $n^{2}=9$ dimensions on the odd side are
the ones supplied by \Cref{thm:rankzero}; the remaining $n(n-1)=6$ come from
the even side and are specific to the line-bundle structure.
\end{corollary}

\begin{proof}
Take $\cE_{1}=\bigoplus L_{i}$ and $\cE_{2}=\bigoplus L_{i}^{-1}$. The index
hypothesis gives $H^{2}(L_{i}^{\pm1}\otimes L_{j}^{\mp1})=0$ for $i\neq j$
and hence the required $\Ext^{2}$ vanishing, and the Chern character
hypothesis is the one of \Cref{thm:rankzero}. The exact decomposition of
the kernel at $n=3$ into an odd part of dimension $9$ and an even part of
dimension $6$ is a computation, recorded in \Cref{sec:verification}: since
$e^{\lambda}$ and $e^{-\lambda}$ combine into $\cosh\lambda$ and
$\sinh\lambda$, which occupy disjoint degrees, the map splits by parity, and
the odd part is exactly the annihilator of \Cref{lem:annihilator} tensored
with the sign vector.
\end{proof}

\begin{remark}[Every object in the family, and what the theorem says]
\label{rem:familyuniform}
The same computation was run over every configuration
$\lambda_{i}=s_{i}\gamma+c_{i}\ell$ with $\sum s_{i}=\sum c_{i}=0$ in a box
of side $7$ whose virtual Chern character lies on the Weil line: eighteen
distinct configurations, every one with kernel exactly $15$. Adding an
$\eta$-component produced no Weil objects at all in a box of side $3$. The
deficiency is not a defect of the particular six line bundles chosen; it is
the same for every member of the family, and \Cref{thm:rankzero} says why.

What the theorem removes is precisely the economy described at the start of
this subsection. An object realising a Weil class through semiregularity
cannot have its Chern character concentrated on the Weil line. It must carry
a nonzero rank, or nonzero Chern character in other degrees, so that
$z\mapsto z\cdot\ch(\cE)$ is injective on $H^{0,2}$; the degree-two component
of that product is $\rk(\cE)\cdot z$, and rank zero is exactly the case in
which it vanishes. Markman's sheaves have nonzero rank, and this is not incidental to their
success: the class he shows to be deformation-invariant is the normalised
character $\exp(-c_{1}(\cE)/\rk\cE)\ch(\cE)$, which divides by the rank,
and the objects at $n=3$ are reflexive sheaves \cite{Mar25a,Mar25b}. So the fourth ingredient of
\Cref{rem:deformation} is narrowed once more: not a difference of two sheaves
whose classes cancel outside the Weil line, but a single sheaf of positive
rank whose Chern character has a Weil component alongside the rest.
\end{remark}

The hypothesis of \Cref{thm:rankzero} is stated for a direct sum because that
is where the element $\xi_{z}$ is visibly nonzero. The mechanism behind it is
more general, and it reaches the object that \Cref{sec:reduction} uses.

\begin{proposition}[Concentrated Chern characters are never semiregular]
\label{prop:concentrated}
Let $A$ be of $(K,-1,n)$-Weil type and let $E$ be a nonzero perfect complex on
$A$ with $\ch_{k}(E)=0$ for $k\ne n$ and $\ch_{n}(E)=N\omega$ for some
$N\ne0$ and some nonzero $\omega\in\HW(A)$. If the natural map
\begin{equation}\label{eq:idmap}
  H^{2}(\cO_{A})\longrightarrow\Ext^{2}(E,E),\qquad z\longmapsto z\cdot\id_{E},
\end{equation}
is injective, then $\dim\ker\sigma_{E}\ge n^{2}$, so $E$ is not semiregular.
The map \eqref{eq:idmap} is injective whenever $E$ is represented by a virtual
sum of line bundles, in particular at every base point of
\Cref{cor:splitalgebraic}, \Cref{thm:quatweil} and \Cref{thm:everyn}.
\end{proposition}

\begin{proof}
For any perfect complex the semiregularity map satisfies
$\sigma_{E}(z\cdot\id_{E})=z\cdot\ch(E)$, since the trace of
$z\cdot\id_{E}\circ\mathrm{At}(E)^{q}$ is $z$ times the degree-$q$ component
of $\ch(E)$. Hence for $z$ in the annihilator of $\omega$ in $H^{0,2}(A)$,
\[
  \sigma_{E}(z\cdot\id_{E})=N\,z\cdot\omega=0 ,
\]
and that annihilator has dimension exactly $n^{2}$ by
\Cref{lem:annihilator}. If \eqref{eq:idmap} is injective the classes
$z\cdot\id_{E}$ are independent, which gives the bound.

For the last sentence, write $E$ in $K_{0}(A)$ as
$\sum_{i}\varepsilon_{i}[L_{i}]$ with $\varepsilon_{i}=\pm1$ and $L_{i}$ line
bundles, realised as a direct sum of the $L_{i}$ and of shifts of the
remaining ones. The map \eqref{eq:idmap} is then $z\mapsto(z,\dots,z)$ into
$\bigoplus_{i}\Ext^{2}(L_{i},L_{i})=\bigoplus_{i}H^{2}(\cO_{A})$, which is
injective. At a split member $\omega_{1}$ is a polynomial in
$\beta,\widehat\beta,\ell$ by \Cref{thm:splitclosed}(iii), at a quaternionic
member it is a polynomial in divisor classes by \Cref{thm:quatweil}, and at a
product it is one by \Cref{thm:everyn}; in each case the Chern character is a
ring expression in classes of line bundles, so such a representation exists.
\end{proof}

\begin{proposition}[The class forces the size of the object]%
\label{prop:objectsize}
Let $A$ be of $(K,-1,n)$-Weil type, let $\omega$ be a nonzero element of the
Weil line and let $E$ be a perfect complex on $A$ with
\[
  \ch(E)=r+N\omega,\qquad r\in H^{0}(A,\QQ),\ N\ne0,
\]
and no other component. Then
\[
  \chi(E,E)\;=\;(-1)^{n}N^{2}\!\int_{A}\omega^{2},
\]
independently of $r$, and consequently
\[
  \sum_{i}\dim\Ext^{i}(E,E)\;\ge\;N^{2}\Bigl|\int_{A}\omega^{2}\Bigr|\;>\;0 .
\]
\end{proposition}

\begin{proof}
The Todd class of an abelian variety is $1$, so Riemann--Roch for the Ext
algebra reads $\chi(E,E)=\int_{A}\ch(E)^{\vee}\ch(E)$ with
$\ch_{k}^{\vee}=(-1)^{k}\ch_{k}$. Expanding,
\[
  \ch(E)^{\vee}\ch(E)
  = r^{2}+\bigl(1+(-1)^{n}\bigr)rN\omega+(-1)^{n}N^{2}\omega^{2},
\]
and of the three terms only the last has degree $4n$, so only it survives the
integral; the rank enters solely through $r^{2}$, which sits in degree zero.
For the inequality, $\chi(E,E)=\sum_{i}(-1)^{i}\dim\Ext^{i}(E,E)$, so the
total dimension is at least $|\chi(E,E)|$. Finally
$\int_{A}\omega^{2}\ne0$: writing $\omega=\alpha_{+}+\alpha_{-}$ with
$\alpha_{\pm}$ spanning $\bigwedge^{2n}V_{\pm}$, each of $\alpha_{\pm}$ is a
top form on its own eigenspace and squares to zero, so
$\omega^{2}=2\alpha_{+}\alpha_{-}$, which is a nonzero multiple of the
fundamental class. The computation is item (XVIII) of
\Cref{sec:verification}.
\end{proof}

\begin{remark}[The constant, and what it is not]\label{rem:sizeconstant}
The bound is on the whole Ext algebra and is insensitive to the rank, which
is the point: enlarging the rank buys no room. The constant
$|\int\omega^{2}|$ depends on the normalisation of $\omega$; measured against
$\eta$ it picks up factors of $d$ and of the polarisation type, so no
dimension-free growth rate should be read off it. What is intrinsic, and what
the bound uses, is only that the Weil class is not isotropic. Together with
\Cref{prop:concentrated} this says that an object carrying the class is
pinned from two sides at once: its semiregularity map has a kernel of
dimension at least $n^{2}$ unless $H^{2}(\cO_{A})$ acts on it degenerately,
and its Ext algebra has dimension at least $N^{2}|\int\omega^{2}|$ whatever
its rank.
\end{remark}

\begin{remark}[What this does to the semiregularity reduction]
\label{rem:killsreduction}
\Cref{thm:semiregclosure} closes the family under the hypothesis that
$\sigma_{E}$ is injective for the complex $E$ of \Cref{prop:ktheory}, whose
Chern character is $N\omega_{1}$ in degree $n$ and zero elsewhere. That is
exactly the hypothesis of \Cref{prop:concentrated}, so at every base point
this paper constructs the hypothesis of \Cref{thm:semiregclosure} fails for
that complex, and fails by the fixed margin $n^{2}$.

Two things keep this from being a refutation of \Cref{thm:semiregclosure}
itself. The hypothesis is about a chosen perfect complex representing a class
in $K_{0}(A)_{\QQ}$, and distinct representatives of one class have distinct
$\Ext^{2}$ groups, so the theorem is not void; what is ruled out is the
representative the construction supplies. And
\Cref{prop:concentrated} needs \eqref{eq:idmap} to be injective, which is a
genuine hypothesis for an object of rank zero: it fails, for instance, for a
sheaf supported at a point. So a representative escaping
\Cref{prop:concentrated} must receive $H^{2}(\cO_{A})$ non-injectively, which
for a rank-zero object means its support is degenerate. We record the
situation as it stands: of the two reductions of \Cref{sec:reduction}, the
degree bound of \Cref{thm:degreebound} is untouched by anything here, and the
semiregularity reduction of \Cref{thm:semiregclosure} is unavailable for every
object we can presently write down.
\end{remark}

\begin{figure}[t]
\centering
  \includegraphics[width=0.94\textwidth]{fig_annihilator}
\caption{The structure behind \Cref{lem:annihilator}, drawn at $n=3$.
$H^{0,2}(A)=\bigwedge^{2}H^{0,1}$ splits into three blocks by how many factors
come from $V_{+}^{0,1}$ and how many from $V_{-}^{0,1}$. The two outer blocks,
of dimension $\binom{n}{2}=3$ each, are carried injectively by multiplication
with the Weil class. The middle block, of dimension $n^{2}=9$ and spanned by
the mixed products, goes to zero, because each such product repeats a factor
of $\alpha_{+}$ and a factor of $\alpha_{-}$. That $n^{2}$ is the deficit in
\Cref{prop:concentrated}.}
\label{fig:annihilator}
\end{figure}


\subsection{The annihilator is the tangent space to the family}
\label{ssec:tangent}

The number $n^{2}$ in \Cref{lem:annihilator} is the dimension of the Weil
family. That is not a coincidence, and this subsection proves it is not: the
polarisation identifies the annihilator with the tangent space to the family,
canonically and in every dimension. Two consequences follow. The first is
that the deficiency of the semiregularity map at an object of rank zero is
not a numerical accident to be repaired by a better choice of object; it is
the family itself, seen inside $H^{0,2}$. The second is that the first-order
Hodge locus of the Weil line is exactly the Weil family, so the directions the
semiregularity map cannot see are precisely the directions along which the
class stays Hodge and would have to be carried.

Write $V_{\pm}$ for the two eigenspaces of $K$ on $H^{1}(A,\CC)$, each of
dimension $2n$, and
\[
  V_{\pm}^{1,0}=V_{\pm}\cap H^{1,0},\qquad
  V_{\pm}^{0,1}=V_{\pm}\cap H^{0,1},
\]
each of dimension $n$ by \Cref{lem:balanced}. Let $E$ denote the
polarisation form on $H^{1}(A,\QQ)$, extended to $H^{1}(A,\CC)$. We use
throughout that $E$ vanishes on $V_{+}\times V_{+}$ and on $V_{-}\times V_{-}$
and pairs $V_{+}$ with $V_{-}$ perfectly, which is condition (iii) of
\Cref{def:weiltype} read on the eigenspaces, and that $E$ vanishes on
$H^{1,0}\times H^{1,0}$ and on $H^{0,1}\times H^{0,1}$, which is the first
Riemann bilinear relation.%
\footnote{The first of these is the computation in the proof of
\Cref{prop:weildomain}: for $x,y\in V_{+}$ one has
$E(\sqrt{-d}\,x,y)=\sqrt{-d}\,E(x,y)$ from the eigenvalue and
$E(\sqrt{-d}\,x,y)=-E(x,\sqrt{-d}\,y)=-\sqrt{-d}\,E(x,y)$ from
\Cref{def:weiltype}(iii), so $E(x,y)=0$; the same on $V_{-}$, and
nondegeneracy of $E$ then forces the pairing between $V_{+}$ and $V_{-}$ to be
perfect. Together the two vanishing statements say that of the four blocks of
$E$ on $(V_{+}^{1,0}\oplus V_{+}^{0,1})\times(V_{-}^{1,0}\oplus V_{-}^{0,1})$
only two are nonzero, namely $V_{+}^{1,0}\times V_{-}^{0,1}$ and
$V_{+}^{0,1}\times V_{-}^{1,0}$, and each of those two is a perfect pairing of
$n$-dimensional spaces. That is the only input to \Cref{thm:anntangent}.}

A first-order deformation of $A$ is a class in $H^{1}(A,T_{A})$, which for an
abelian variety is canonically $\Hom(H^{1,0},H^{0,1})$, since the tangent
bundle is trivial and $H^{1,0}$ is the dual of its space of sections. Such a
$v$ is extended to $H^{1}(A,\CC)$ by zero on $H^{0,1}$ and acts on
$H^{k}(A,\CC)=\bigwedge^{k}H^{1}(A,\CC)$ as the derivation it generates; we
write $v\lrcorner\,\xi$ for that action. A Hodge class $\xi$ of type $(p,p)$
remains of type $(p,p)$ to first order along $v$ exactly when
$v\lrcorner\,\xi=0$. We call $v$ \emph{polarised} when
$E(vx,y)+E(x,vy)=0$ for all $x,y\in H^{1,0}$, which is the condition that the
deformation preserve the polarisation, and the polarised $v$ form a space of
dimension $n(2n+1)=\dim\cA_{2n}$.

\begin{proposition}[The first-order Hodge locus of the Weil line]
\label{prop:firstorder}
Let $A$ be of $(K,-1,n)$-Weil type and let $v$ be a polarised first-order
deformation of $A$. The following are equivalent.
\begin{enumerate}[label=\textup{(\roman*)},nosep]
\item $v\lrcorner\,\omega=0$ for every $\omega$ in the Weil line.
\item $v$ commutes with the action of $K$.
\item $v\bigl(V_{+}^{1,0}\bigr)\subseteq V_{+}^{0,1}$ and
  $v\bigl(V_{-}^{1,0}\bigr)\subseteq V_{-}^{0,1}$.
\end{enumerate}
The space of such $v$ has dimension $n^{2}$, and under
\Cref{prop:weildomain} it is the tangent space at $[A]$ to the Weil family
$\cD_{n,n}$. Already one of the two rational generators of the Weil line cuts
it out.
\end{proposition}

\begin{proof}
The equivalence of (ii) and (iii) is immediate: $K$ acts on $V_{\pm}$ by the
two scalars $\pm\sqrt{-d}$, so commuting with $K$ is preserving $V_{+}$ and
$V_{-}$, and the extension of $v$ kills $H^{0,1}$, so the content is where it
sends $V_{\pm}^{1,0}$.

For (i) and (iii), let $\alpha_{+}$ generate $\bigwedge^{2n}V_{+}$ and write
$\alpha_{+}=e_{1}\wedge\cdots\wedge e_{2n}$ for a basis $e_{1},\dots,e_{2n}$
of $V_{+}$ adapted to $V_{+}=V_{+}^{1,0}\oplus V_{+}^{0,1}$. As a derivation,
\[
  v\lrcorner\,\alpha_{+}
  =\sum_{j=1}^{2n}e_{1}\wedge\cdots\wedge v(e_{j})\wedge\cdots\wedge e_{2n}.
\]
Split $v(e_{j})$ along $H^{1}(A,\CC)=V_{+}\oplus V_{-}$. The $V_{+}$
components contribute
$\operatorname{tr}\bigl(\pi_{+}\circ v|_{V_{+}}\bigr)\,\alpha_{+}$, and that
trace vanishes: the map $\pi_{+}\circ v|_{V_{+}}$ kills $V_{+}^{0,1}$ and
carries $V_{+}^{1,0}$ into $V_{+}^{0,1}$, so in the adapted basis it is
strictly block off-diagonal and has zero diagonal. Writing
$v(e_{j})_{-}=\sum_{k}c_{jk}f_{k}$ for a basis $f_{1},\dots,f_{2n}$ of
$V_{-}$, the $V_{-}$ components contribute
\[
  \sum_{j,k}\pm\,c_{jk}\;
  e_{1}\wedge\cdots\wedge\widehat{e_{j}}\wedge\cdots\wedge e_{2n}\wedge f_{k},
\]
and the $4n^{2}$ monomials appearing here are distinct members of the
standard basis of $\bigwedge^{2n}H^{1}(A,\CC)$, hence linearly independent.
So $v\lrcorner\,\alpha_{+}=0$ if and only if every $c_{jk}$ vanishes, which is
$v(V_{+})\subseteq V_{+}$. The same argument on $\alpha_{-}$ gives
$v\lrcorner\,\alpha_{-}=0$ if and only if $v(V_{-})\subseteq V_{-}$. The Weil
line spans $\bigwedge^{2n}V_{+}\oplus\bigwedge^{2n}V_{-}$ over $\CC$ by
\Cref{prop:weilline}, and each of the two rational generators of
\Cref{thm:explicitgens} has both components nonzero, so a single generator
already forces both conditions. This is (i) if and only if (iii).

For the dimension, let $v$ satisfy (iii) and write $v_{+}$ and $v_{-}$ for
its two blocks. If $x,y$ both lie in $V_{+}^{1,0}$ then $v_{+}x$ and $y$ both
lie in $V_{+}$, so both terms of the polarised condition vanish, and likewise
if both lie in $V_{-}^{1,0}$. For $x\in V_{+}^{1,0}$ and $y\in V_{-}^{1,0}$
the condition reads $E(v_{+}x,y)=-E(x,v_{-}y)$. The left side is determined by
$v_{+}$, and the pairing between $V_{+}^{1,0}$ and $V_{-}^{0,1}$ is perfect,
so this determines $v_{-}y$ for each $y$ and hence determines $v_{-}$ from
$v_{+}$. So the polarised deformations satisfying (iii) are a space of
dimension $\dim\Hom(V_{+}^{1,0},V_{+}^{0,1})=n^{2}$. By
\Cref{prop:weildomain} the family $\cD_{n,n}$ is an open subset of the
Grassmannian of $n$-planes in $V_{+}$, whose tangent space at
$V_{+}^{1,0}$ is $\Hom(V_{+}^{1,0},V_{+}/V_{+}^{1,0})=\Hom(V_{+}^{1,0},V_{+}^{0,1})$,
and the identification just made is that isomorphism.
\end{proof}

\begin{proposition}[Infinitesimal rigidity of divisor classes]%
\label{prop:rigidity}
Keep the notation above and write a class $\delta\in H^{1,1}$ in the four
blocks
\[
  a\ \text{on}\ P\otimes\bbar{P},\quad
  b\ \text{on}\ P'\otimes\bbar{P}',\quad
  c\ \text{on}\ P\otimes\bbar{P}',\quad
  e\ \text{on}\ P'\otimes\bbar{P},
\]
the first two carrying the norm character of $\Kx$ and the last two the
characters $\tau^{2}$ and $\bbar\tau^{\,2}$; reality makes $a$ and $b$
anti-hermitian and ties $e$ to $-\bbar{c}^{\,T}$. Identify the tangent space
to $\cD_{n,n}$ with the $n\times n$ matrices $V$ as in
\Cref{thm:anntangent}. Then contraction with a tangent vector is
\begin{equation}\label{eq:phidelta}
  \Phi_{\delta}(V)
  \;=\;\bigl(\,Va+b^{T}V,\ \ \operatorname{alt}(Vc),\ \
       \operatorname{alt}(-V^{T}e)\,\bigr),
\end{equation}
with values in
$\bbar{P}'\otimes\bbar{P}\oplus\bigwedge^{2}\bbar{P}'\oplus\bigwedge^{2}\bbar{P}
=H^{0,2}$, and $\ker\Phi_{\delta}$ is the tangent space to the locus where
$\delta$ stays of type $(1,1)$. Consequently:
\begin{enumerate}[label=\textup{(\roman*)},nosep]
\item $\Phi_{\delta}=0$ if and only if $\delta$ is a real multiple of $\eta$.
  The rigid classes form a real line, which is the infinitesimal form of the
  statement that $\NS(A)_{\QQ}$ has rank one at the general member.
\item For the norm-character class with block $a$ and $b=c=e=0$,
  \[
    \dim\ker\Phi_{\delta}\;=\;n\,(n-\operatorname{rank}a),
  \]
  so such a class cuts a locus of codimension $n\operatorname{rank}a$, at
  least $n$.
\item A general class has $\ker\Phi_{\delta}=0$: it rigidifies the family to
  a point.
\end{enumerate}
\end{proposition}

\begin{proof}
For \eqref{eq:phidelta}, a polarised deformation commuting with $K$ sends $P$
into $\bbar{P}'$ and $P'$ into $\bbar{P}$, and the symmetry
$E(vx,y)=E(x,vy)$ forces the second matrix to be $-V^{T}$, which is
\Cref{thm:anntangent}. Applying $v$ to the $(1,0)$ slot of each block and
collecting terms in $H^{0,2}$ gives the three displayed components; the sign
in the first comes from $\bbar{P}\wedge\bbar{P}'=-\bbar{P}'\wedge\bbar{P}$.

(i) If $\Phi_{\delta}=0$ then $\operatorname{alt}(Vc)=0$ for every $V$, which
for $n\ge2$ forces $c=0$, and likewise $e=0$. The first component gives
$Va=-b^{T}V$ for every $V$; taking $V$ to run over the matrix units shows
$a$ is a scalar $\lambda$ and $b^{T}=-\lambda$. Anti-hermitian scalars are
the purely imaginary ones, so the solutions form a real line. It contains
$\eta$, whose blocks are $a=i$, $b=-i$, and the two therefore agree.

(ii) With $b=c=e=0$ the condition is $Va=0$, that is, the rows of $V$
annihilate the column space of $a$; the solution space has dimension
$n(n-\operatorname{rank}a)$.

(iii) Immediate from (i) and (ii), and checked on samples.

All three, together with the value $n^{2}$ for the dimension of the tangent
space, are verified over $\QQ(i)$ in exact arithmetic for $n=2,3,4$ in item
(XIX) of \Cref{sec:verification}; at $n=4$ the kernels in (ii) are $12$, $8$,
$4$, $0$ for ranks $1$ to $4$.
\end{proof}

\begin{remark}[Why the count matters here]\label{rem:rigiditywhy}
\Cref{prop:rigidity}(ii) is the reason the divisorial supply of
\Cref{thm:divcriterion} cannot be enlarged by a better choice of class. Any
norm-character class other than a multiple of $\eta$ has
$\operatorname{rank}a\ge1$ and so cuts codimension at least $n$; the locus it
defines is proper, and by (i) no class except the polarisation avoids cutting
at all. So each new class buys a subvariety, never an open set, which is
\Cref{prop:divfails} again with a dimension attached.
\end{remark}

\begin{figure}[htbp]
\centering
  \includegraphics[width=0.46\textwidth]{fig_fibres}
\caption{The geometry behind \Cref{lem:annihilator}, drawn for the real shadow of
the hermitian form. The ruled surface is the null cone of $H$; the two
transverse planes are the conjugate isotropic subspaces $V_{+}$ and $V_{-}$,
which conjugation interchanges, so neither is defined over $\QQ$ and only
their sum is. The segments are the fibres joining a vector of one to its
conjugate in the other, and the strip they sweep is the mixed block
$V_{+}^{0,1}\otimes V_{-}^{0,1}$, of dimension $n^{2}$, which multiplication
by the Weil class sends to zero and which is the deficit in
\Cref{prop:concentrated}. \protect\footnotemark}
\label{fig:fibres}
\end{figure}
\footnotetext{The cone is drawn with its rulings so that it reads as the ruled surface
it is rather than as a shaded volume. The two planes meet only at the origin,
so the fibres are disjoint away from it and the strip they sweep is a cone;
that is why the block is a tensor product and not a direct sum.}

\begin{theorem}[The annihilator is the tangent space]\label{thm:anntangent}
Let $A$ be of $(K,-1,n)$-Weil type and let $\omega$ be a rational generator of
the Weil line. The polarisation induces an isomorphism
\[
  T_{[A]}\cD_{n,n}\;\xrightarrow{\ \sim\ }\;
  \operatorname{Ann}_{H^{0,2}(A)}(\omega),
\]
canonical in the sense that it depends on no choice of basis and on no choice
of $\omega$ on the line. Explicitly, if $p_{1},\dots,p_{n}$ is a basis of
$V_{+}^{1,0}$ and $q_{1},\dots,q_{n}$ the basis of $V_{-}^{0,1}$ dual to it
under $E$, then $\varphi\in\Hom(V_{+}^{1,0},V_{+}^{0,1})$ corresponds to
\[
  \sum_{a=1}^{n}q_{a}\wedge\varphi(p_{a})\;\in\;
  V_{-}^{0,1}\otimes V_{+}^{0,1}\;\subset\;\textstyle\bigwedge^{2}H^{0,1}
  =H^{0,2}(A).
\]
\end{theorem}

\begin{proof}
The form $E$ pairs $V_{+}$ with $V_{-}$ perfectly and vanishes on
$V_{+}^{1,0}\times V_{-}^{1,0}$, the latter because both spaces lie in
$H^{1,0}$. Hence $E$ restricted to $V_{+}^{1,0}\times V_{-}$ kills
$V_{-}^{1,0}$ and descends to a pairing
\[
  V_{+}^{1,0}\times V_{-}^{0,1}\longrightarrow\CC ,
\]
which is perfect, both spaces having dimension $n$ and the total pairing being
nondegenerate. This gives a canonical isomorphism
$\bigl(V_{+}^{1,0}\bigr)^{*}\cong V_{-}^{0,1}$, and therefore
\[
  T_{[A]}\cD_{n,n}
  =\Hom\bigl(V_{+}^{1,0},V_{+}^{0,1}\bigr)
  =\bigl(V_{+}^{1,0}\bigr)^{*}\otimes V_{+}^{0,1}
  \;\cong\;V_{-}^{0,1}\otimes V_{+}^{0,1},
\]
the first equality being the tangent space computation at the end of the proof
of \Cref{prop:firstorder}. Since $H^{0,1}=V_{+}^{0,1}\oplus V_{-}^{0,1}$, the
map $q\otimes u\mapsto q\wedge u$ identifies $V_{-}^{0,1}\otimes V_{+}^{0,1}$
with the middle summand of
$\bigwedge^{2}H^{0,1}$, and that summand is the annihilator by
\Cref{lem:annihilator}. The displayed formula is the composite of these
isomorphisms written in a dual pair of bases, so it is independent of the
basis chosen; and \Cref{lem:annihilator} gives the same subspace for every
$\omega$ on the line, so the isomorphism does not depend on $\omega$ either.
\end{proof}

\begin{figure}[tbp]
\centering
  \includegraphics[width=\textwidth]{fig_tangent}
\caption{\Cref{thm:anntangent}. The polarisation has only two nonzero blocks
on $V_{+}\times V_{-}$: it pairs $V_{+}^{1,0}$ with $V_{-}^{0,1}$ and
$V_{+}^{0,1}$ with $V_{-}^{1,0}$, each perfectly, and the other two blocks
vanish by the first Riemann relation. The upper right block turns the tangent
space $\Hom(V_{+}^{1,0},V_{+}^{0,1})$ to the Weil family into
$V_{-}^{0,1}\otimes V_{+}^{0,1}$, which is the middle block of
$\bigwedge^{2}H^{0,1}$ drawn in \Cref{fig:annihilator} and the annihilator of
the Weil line. The two $n^{2}$ dimensional spaces are therefore the same
subspace of $H^{0,2}$, so the directions the semiregularity map cannot see are
the directions of the family along which the class stays Hodge.}
\label{fig:tangent}
\end{figure}


\begin{corollary}[The kernel contains the tangent space to the family]
\label{cor:kernelfamily}
Let $A$ be of $(K,-1,n)$-Weil type and let $E$ be a nonzero perfect complex on
$A$ satisfying the hypotheses of \Cref{prop:concentrated}, that is, with Chern
character concentrated on the Weil line and with \eqref{eq:idmap} injective.
Then $\ker\sigma_{E}$ contains a subspace canonically isomorphic to
$T_{[A]}\cD_{n,n}$, namely the image of the composite
\[
  T_{[A]}\cD_{n,n}\;\cong\;\operatorname{Ann}_{H^{0,2}}(\omega)
  \;\hookrightarrow\;H^{2}(\cO_{A})\;\xrightarrow{\;z\mapsto z\cdot\id_{E}\;}\;
  \Ext^{2}(E,E).
\]
By \Cref{prop:firstorder} that subspace is also the tangent space to the
first-order Hodge locus of the class $\ch_{n}(E)$.
\end{corollary}

\begin{proof}
Immediate from \Cref{thm:anntangent}, \Cref{prop:concentrated} and
\Cref{prop:firstorder}.
\end{proof}

\begin{theorem}[The weakened criterion is unavailable for concentrated
Chern characters]\label{thm:weakfails}
Let $A$ be of $(K,-1,n)$-Weil type and let $E$ be a nonzero perfect complex on
$A$ satisfying the hypotheses of \Cref{prop:concentrated}. Let
\[
  \operatorname{ev}_{E}\colon HH^{2}(A)\longrightarrow\Ext^{2}(E,E)
\]
be the evaluation on $E$ of the natural transformations $\id\to\id[2]$ of
$D^{b}(A)$, as in \cite[Section 11.4]{Mar25b}. Then
\[
  \operatorname{ev}_{E}\bigl(H^{2}(\cO_{A})\bigr)\cap\ker\sigma_{E}
  =\Delta:=\bigl\{\,z\cdot\id_{E}\;:\;z\in\Ann_{H^{0,2}}(\omega)\,\bigr\},
\]
a subspace of dimension $n^{2}$. In particular $\sigma_{E}$ is not injective
on the image of $\operatorname{ev}_{E}$, so $E$ satisfies neither the
hypothesis of \cite[Lemma 11.3]{Mar25b} nor the weakened hypothesis of
\cite[Question 11.4]{Mar25b}.
\end{theorem}

\begin{proof}
The Hochschild cohomology $HH^{*}(A)=\Ext^{*}_{A\times A}(\cO_{\Delta},\cO_{\Delta})$
is the algebra of natural transformations of the identity functor, and the
Hochschild--Kostant--Rosenberg isomorphism \cite{Cal05,BF08} writes
$HH^{2}(A)\cong H^{2}(\cO_{A})\oplus H^{1}(A,T_{A})\oplus H^{0}(A,\bigwedge^{2}T_{A})$,
with no correction from the Todd class because $T_{A}$ is trivial. The first
summand is the image of $H^{2}(\cO_{A})=\Ext^{2}(\cO_{A},\cO_{A})$ under
pullback to $A\times A$, and the natural transformation it defines sends a
complex $F$ to the composite $F=F\otimes\cO_{A}\xrightarrow{\;\id\otimes z\;}
F\otimes\cO_{A}[2]=F[2]$, that is to $z\cdot\id_{F}$; this is the statement
that the action of $HH^{*}(A)$ on $\Ext^{*}(F,F)$ restricts on $H^{*}(\cO_{A})$
to the $\cO_{A}$-module structure.%
\footnote{Concretely: a class $z\in H^{2}(\cO_{A})$ is represented by a
$2$-extension of $\cO_{A}$ by itself, its pullback $\pi_{1}^{*}z$ to
$A\times A$ is a $2$-extension of $\cO_{A\times A}$ by itself, and tensoring
with $\cO_{\Delta}$ gives a $2$-extension of $\cO_{\Delta}$ by itself, which
is the element of $HH^{2}(A)$ in question. The Fourier--Mukai transform with
kernel $\cO_{\Delta}$ is the identity, and the transform with kernel the
extension is the extension of the identity by itself twisted by $z$, so its
value on $F$ is $z\cdot\id_{F}$. Markman uses exactly this identification when
he writes the left vertical arrow of the diagram in the proof of his Lemma
11.3 as HKR followed by contraction with $\ch(E)$: on the summand $H^{2}(\cO_{A})$ contraction is the
cup product $z\mapsto z\cdot\ch(E)$, which is
$\sigma_{E}(z\cdot\id_{E})$ by the computation in the proof of
\Cref{thm:rankzero}.}
Hence $\operatorname{ev}_{E}(z)=z\cdot\id_{E}$ for $z\in H^{2}(\cO_{A})$, and
$\sigma_{E}(z\cdot\id_{E})=z\cdot\ch(E)=N\,z\cdot\omega$ by the proof of
\Cref{prop:concentrated}, which vanishes exactly for
$z\in\Ann_{H^{0,2}}(\omega)$. So
$\operatorname{ev}_{E}(H^{2}(\cO_{A}))\cap\ker\sigma_{E}=\Delta$, and
$\dim\Delta=\dim\Ann_{H^{0,2}}(\omega)=n^{2}$ by \Cref{lem:annihilator} and
the injectivity of \eqref{eq:idmap}. Since $\Delta\ne0$, the restriction of
$\sigma_{E}$ to the image of $\operatorname{ev}_{E}$ has a nonzero kernel.
The hypothesis of \cite[Lemma 11.3]{Mar25b} is that
$\ker\operatorname{ev}_{E}$ equals the kernel of contraction with $\ch(E)$ and
that $\operatorname{ev}_{E}$ is surjective; the second condition would make
$\sigma_{E}$ injective on all of $\Ext^{2}(E,E)$, contradicting
\Cref{prop:concentrated}, and the weakened hypothesis of
\cite[Question 11.4]{Mar25b} is injectivity of $\sigma_{E}$ on the image,
which has just been shown to fail.
\end{proof}

\begin{theorem}[A line bundle summand off the polarisation line obstructs]
\label{thm:linebundle}
Let $A$ be of $(K,-1,n)$-Weil type with $n\ge2$ and let $T=T_{[A]}\cD_{n,n}$.
Let $E$ be a perfect complex on $A$ with a direct summand $L[p]$, $L$ a line
bundle, and suppose that $\ch(E)$ remains of Hodge type to first order along
$\cD_{n,n}$, that is $v\lrcorner\,\ch(E)=0$ for every $v\in T$. If $c_{1}(L)$
is not a real multiple of $\eta$, then there is $v\in T$ with
\[
  \operatorname{ev}_{E}(v)\ne0
  \qquad\text{and}\qquad
  \sigma_{E}\bigl(\operatorname{ev}_{E}(v)\bigr)=0 .
\]
Consequently $E$ is not semiregular, does not satisfy the weakened hypothesis
of \cite[Question 11.4]{Mar25b}, and is obstructed to first order along $v$.
The same holds for a twisted complex, with $c_{1}(L)$ replaced by the twisted
first Chern class of the summand, which is $c_{1}(L)$ shifted by the class
common to all summands.
\end{theorem}

\begin{proof}
The class $c_{1}(L)$ is a real class of type $(1,1)$ and not a real multiple
of $\eta$, so by \Cref{prop:rigidity}(i) the contraction
$\Phi_{c_{1}(L)}\colon T\to H^{0,2}(A)$ is not zero: there is $v\in T$ with
$v\lrcorner\,c_{1}(L)\ne0$. On the summand $H^{1}(A,T_{A})$ of $HH^{2}(A)$ the
evaluation map is contraction with the Atiyah class,
$\operatorname{ev}_{E}(v)=v\lrcorner\,\mathrm{At}(E)$; this is the top row of
the diagram of Buchweitz and Flenner reproduced as \cite[(2.1)]{Mar25b}, and
its identification with the Hochschild evaluation is \cite{BF08}, used by
Markman in the proof of his Lemma 11.3. The Atiyah class of a direct sum is
the direct sum of the Atiyah classes, and that of $L[p]$ is $c_{1}(L)$, so the
component of $\operatorname{ev}_{E}(v)$ in the summand
$\Ext^{2}(L[p],L[p])=H^{0,2}(A)$ of $\Ext^{2}(E,E)$ is
$v\lrcorner\,c_{1}(L)\ne0$. On the other hand the same diagram gives
$\sigma_{E}(\operatorname{ev}_{E}(v))=v\lrcorner\,\ch(E)$, which vanishes by
hypothesis. So $\operatorname{ev}_{E}(v)$ is a nonzero element of
$\ker\sigma_{E}\cap\image\operatorname{ev}_{E}$; this is the failure of both
hypotheses, and $v\lrcorner\,\mathrm{At}(E)\ne0$ is the statement that $E$
does not lift to a first-order deformation of the pair along $v$, as recalled
in \cite[Section 2]{Mar25b}. Twisting by a class $B$ replaces the Atiyah
class of every rank one summand by its shift by $B$ and $\ch(E)$ by the
twisted Chern character, and nothing else in the argument changes.
\end{proof}

\begin{remark}[What this measures, and what it does not settle]
\label{rem:whatmeasures}
The semiregularity criterion converts a statement about the Chern character
into a statement about the object, and it does so by requiring
$\sigma_{E}$ to be injective. \Cref{cor:kernelfamily} says that for an object whose
Chern character is concentrated on the Weil line the requirement fails, and
that the failure is measured by the tangent space to the very family across
which the class is to be transported, in every dimension and with no
hypothesis beyond \eqref{eq:idmap}. That is a statement about the size and the
position of the kernel; by itself it is not a statement that the obstruction
classes are nonzero, and it does not show that the Weil class fails to
deform. For the explicit object of \Cref{prop:explicitobject} the obstruction
classes are computed in \Cref{prop:splitobstruction}: they are nonzero in
every direction of the family normal to the split locus, they lie in
$\ker\sigma_{E}$, and they meet the subspace of \Cref{cor:kernelfamily} only
in zero. The Weil class deforms along the family, by \Cref{prop:firstorder};
that object does not.

Nor does it decide the weaker criterion in the form that matters. Markman
asks, in \cite[Question 11.4]{Mar25b}, whether the semiregularity theorem
survives when $\sigma_{E}$ is required to be injective only on the image of
the evaluation map $\operatorname{ev}_{E}$ rather than on all of
$\Ext^{2}(E,E)$, and records in his Section 12 that an affirmative answer
would carry the method to higher dimensions and to fields of larger degree.
\Cref{thm:weakfails} says that the classes of \Cref{cor:kernelfamily} lie in
that image, so for concentrated Chern characters the weakened hypothesis
fails as well, in every dimension, by the same $n^{2}$ dimensions; and
\Cref{thm:linebundle} says that it fails for every object with a line bundle
summand off the polarisation line, whatever the rank, at every member of the
family, which by \Cref{cor:noline} removes every sum of line bundles from
consideration at once. For the
explicit split object of \Cref{prop:explicitobject} the failure is complete:
\Cref{prop:evaluation} shows that every one of the fifteen dimensions of
$\ker\sigma_{E}$ is the value of $\operatorname{ev}_{E}$ on a Hochschild
class that already preserves $\ch(E)$ to first order. What is not decided,
and what the question is really about, is whether an object of positive rank
can satisfy the weakened hypothesis without satisfying the original one.
Markman's own sheaves are of positive rank and are not of the concentrated
shape, which is \Cref{rem:familyuniform}; he avoids the question for them by
descending to a quotient on which $\sigma_{E}$ is injective on an invariant
subspace containing the image of $\operatorname{ev}_{E}$ \cite[Proposition
11.5]{Mar25b}, and he proposes in his Section 12 a rank one complex on a
principally polarised abelian fourfold, built from $(d+9)/2$ translates of the
theta divisor, as a candidate for the weakened criterion in dimension eight,
with the verification still to be done. \Cref{thm:candidatefails} carries
out that verification and finds that the candidate fails, for every
discriminant, at every isolated double point of its support.

Voisin's account of the present state of the subject \cite{Voi26} makes the
same point qualitatively and for the variational Hodge conjecture at large,
that semiregular objects are hard to construct and too restricted to carry a
general argument. \Cref{cor:kernelfamily} is that observation in quantitative
form for one family of classes: the restriction is exactly $n^{2}$
dimensions, and those dimensions are the family. Totaro's account
\cite{Tot26} is more optimistic about pushing Markman's method further, and
records that it needed the semiregularity theorem for twisted sheaves, which
Pridham supplied \cite{Pri24}.
\end{remark}

